\begin{frame}{例子：这门课最后到底希望学生带走什么}
  \Core{课程总结不应只回顾目录，而应明确学生已经获得的一套可重复的方法。}
  \begin{block}{不是记住若干论文名字}
    论文只是样本，真正带走的应是分析系统问题的框架。
  \end{block}
  \pause
  \begin{block}{也不是掌握一套固定 runtime}
    runtime 会变，但对象、控制点、收益链路和验证方法这套框架可以迁移。
  \end{block}
  \pause
  \begin{block}{课程结论}
    这门课真正的产出，是一套进入真实推理系统后还能继续使用的工作方法。
  \end{block}
\end{frame}

\begin{frame}{例子：一个优秀最终汇报通常长什么样}
  \Core{好的最终汇报通常不追求面面俱到，而是把问题和证据压到很实。}
  \begin{block}{开场}
    先用一句话讲清系统对象、控制点和目标指标。
  \end{block}
  \pause
  \begin{block}{主体}
    用最小实现和最关键实验支撑核心因果链，而不是堆大量细碎结果。
  \end{block}
  \pause
  \begin{block}{结尾}
    诚实说明失败条件、边界和后续路线，往往比强行夸大收益更有研究价值。
  \end{block}
\end{frame}

\begin{frame}{例子：一个看似很努力但不合格的汇报}
  \Core{课程总结课也要明确失败模式，这样学生才知道之后如何避免。}
  \begin{block}{问题一}
    讲了很多实现细节，却没有说明到底改变了哪个控制点。
  \end{block}
  \pause
  \begin{block}{问题二}
    图很多，但 workload、对照和失败边界都不清楚。
  \end{block}
  \pause
  \begin{block}{问题三}
    把工程量直接当成研究贡献，导致整场报告没有中心判断。
  \end{block}
\end{frame}

\begin{frame}{代码例子：最终汇报推荐结构}
  \Core{课程最后一讲可以给出一份较硬的汇报骨架，帮助学生把材料组织成研究叙事。}
  \begin{block}{结构模板}
    {\ttfamily\small
    \detokenize{1. Problem and system object}\par
    \detokenize{2. Changed control point / mechanism}\par
    \detokenize{3. Minimal implementation path}\par
    \detokenize{4. Evidence, boundary, next step}\par
    }
  \end{block}
  \pause
  \begin{block}{课堂应追问}
    为什么这个结构里没有单独的“工程量展示”章节？
  \end{block}
  \pause
  \begin{block}{课程结论}
    汇报的中心是研究判断，而不是开发过程流水账。
  \end{block}
\end{frame}

\begin{frame}{代码例子：课程项目资料如何互相对齐}
  \Core{你前面要求的“讲义、tutorial、实验和代码互相呼应”，在最后汇报里必须显式体现。}
  \begin{block}{对齐关系}
    {\ttfamily\small
    \detokenize{slides -> problem framing}\par
    \detokenize{tutorials -> concept and code map}\par
    \detokenize{experiments -> evidence and procedure}\par
    \detokenize{code -> control point and implementation}\par
    }
  \end{block}
  \pause
  \begin{block}{课堂应追问}
    如果你的最终报告无法指出这四层材料各自承担什么角色，说明前期材料组织仍然是断裂的。
  \end{block}
  \pause
  \begin{block}{教学意义}
    课程闭环不是文档美观问题，而是研究训练是否真正闭合的问题。
  \end{block}
\end{frame}

\begin{frame}{例子：从课程项目继续做研究的三条路线}
  \Core{课程结束后，真正有潜力的项目通常沿着问题深化、系统加深和证据扩展三条路线继续走。}
  \begin{block}{路线一：问题深化}
    把一个课程级现象收紧成更明确的系统假设。
  \end{block}
  \pause
  \begin{block}{路线二：系统加深}
    从教学代码进入真实 runtime 或异构平台，验证机制是否仍成立。
  \end{block}
  \pause
  \begin{block}{路线三：证据扩展}
    把最小实验扩大到更真实 workload、更严格对照和更强失败分析。
  \end{block}
\end{frame}

\begin{frame}{课堂练习：给未来的自己留一张研究路线卡}
  \Core{课程最后的练习，不是再做一题，而是把当前项目转成可继续推进的路线图。}
  \begin{block}{练习一}
    写出你项目里最值得继续研究的一条假设。
  \end{block}
  \pause
  \begin{block}{练习二}
    写出下一步最值得补的一个实验和一个代码入口。
  \end{block}
  \pause
  \begin{block}{练习三}
    用一句话回答：这门课之后，你进入一个新推理系统时会先看什么？
  \end{block}
\end{frame}
